\pagestyle{empty}
\tight
\begin{tblr}{width=\textwidth,colspec={|Q[c,m,1.9cm]|X[l,m]|},hlines,vlines,hspan=minimal,row{1}={bg=headblue},rowsep=1pt,colsep=3pt}
\SetCell{c,m}\hfil\logo\hfil & \textbf{POLITEKNIK NEGERI MALANG}\par\textbf{JURUSAN TEKNOLOGI INFORMASI}\par\textbf{PROGRAM STUDI: D4 TEKNIK INFORMATIKA} \\
\SetCell[c=2]{c} \textbf{RENCANA PEMBELAJARAN SEMESTER (RPS)} & \\
\end{tblr}
\begin{longtblr}{width=\textwidth,colspec={|Q[l,m,1.9cm]|X[0.85,l,m]|X[1.45,l,m]|X[0.75,l,m]|X[1.15,l,m]|},hlines,vlines,hspan=minimal,rowsep=1pt,colsep=2pt,row{1,3}={bg=headgray,font=\bfseries}}
MATA KULIAH & KODE & BOBOT (SKS)/jam & SEMESTER & TGL. PENYUSUNAN \\
Bahasa Inggris Persiapan Kerja & RTI258002 & 2 SKS/4 jam & 8 & 1 Juli 2025 \\
OTORISASI & Dosen Pengembang RPS & Koordinator RMK & \SetCell[c=2]{l} Ka PRODI &  \\
 & Farid Angga Pribadi, S.Kom., M.Kom. &  & \SetCell[c=2]{l}{\vspace{8mm}\par Dr. Ely Setyo Astuti, S.T., M.T.} &  \\
\SetCell[r=9]{l} Capaian Pembelajaran (CP) & \SetCell[c=4]{l,bg=headgray,font=\bfseries} Capaian Pembelajaran Lulusan Program Studi (CPL-Prodi) &  &  &  \\
 & CPL01 & \SetCell[c=3]{l} Mampu menerapkan nilai ketakwaan, kesadaran hukum, kedisiplinan, profesionalisme, etika profesi, dan pembelajaran sepanjang hayat, serta tanggap terhadap isu sosial dan perkembangan teknologi. &  &  \\
 & CPL03 & \SetCell[c=3]{l} Mampu mengelola tim dan diri sendiri secara profesional serta mengomunikasikan dan mempresentasikan gagasan maupun hasil kerja secara efektif baik lisan maupun tulisan. &  &  \\
 & \SetCell[c=4]{l,bg=headgray,font=\bfseries} Capaian Pembelajaran mata kuliah (CPL-MK) &  &  &  \\
 & CPMK0303 & \SetCell[c=3]{l} Mahasiswa mampu mengomunikasikan dan mempresentasikan hasil kerja teknologi informasi secara efektif dalam Bahasa Inggris baik lisan maupun tulisan. &  &  \\
 & CPMK0102 & \SetCell[c=3]{l} Mahasiswa mampu menerapkan etika profesi, profesionalisme, dan komunikasi lintas budaya dalam lingkungan kerja TI global. &  &  \\
 & \SetCell[c=4]{l,bg=headgray,font=\bfseries} Kemampuan akhir tiap tahapan belajar (Sub-CPMK) &  &  &  \\
 & SCPMK0303-05901 & \SetCell[c=3]{l} Mahasiswa mampu menyusun dokumen profesional dalam Bahasa Inggris meliputi CV, cover letter, portofolio, dan email profesional. &  &  \\
 & SCPMK0102-05902 & \SetCell[c=3]{l} Mahasiswa mampu melakukan simulasi wawancara kerja, presentasi teknis, dan komunikasi dalam tim multikultural secara efektif dalam Bahasa Inggris. &  &  \\
\SetCell[r=2]{l} Penilaian & \SetCell[c=4]{l} *Nilai CPMK didapat dari agregasi nilai SUB-CPMK yang dibebankan &  &  &  \\
 & \SetCell[c=4]{l}{\tiny\begin{tblr}{width=\linewidth,colspec={|Q[c,m,0.1\linewidth]|Q[c,m,0.16\linewidth]|X[l,m]|X[l,m]|X[l,m]|X[l,m]|X[l,m]|X[l,m]|Q[c,m,0.08\linewidth]|},hlines,vlines,hspan=minimal,rowsep=1pt,colsep=0.7pt,row{1,2}={bg=headgray,font=\bfseries}}
Kode CPMK & Kode SCPMK & \SetCell[c=6]{c} Bobot per Bentuk Penilaian & & & & & & Total Bobot \\
 & & Tugas Dokumen & Kuis 1 & UTS & Kuis 2 & PBL Mock & UAS & \\
CPMK0303 & SCPMK0303-05901 & 18\%: dokumen profesional & 5\%: CV dan cover letter & 8\%: portofolio awal & & 10\%: presentasi teknis & 10\%: final portofolio & 51\% \\
CPMK0102 & SCPMK0102-05902 & 12\%: komunikasi profesi & & 7\%: komunikasi kerja & 5\%: interview skills & 12\%: simulasi wawancara & 13\%: final mock interview & 49\% \\
\SetCell[c=8]{r}\textbf{Total Per Penilaian} & & & & & & & & \textbf{100\%} \\
\end{tblr}} &  &  &  \\
Diskripsi Singkat Mata Kuliah & \SetCell[c=4]{l} Bahasa Inggris Persiapan Kerja mempersiapkan mahasiswa memasuki pasar kerja TI global dengan membangun kemampuan komunikasi profesional dalam Bahasa Inggris mencakup penulisan dokumen karier, presentasi teknis, simulasi wawancara, dan komunikasi lintas budaya. &  &  &  \\
Bahan Kajian / Materi Pembelajaran & \SetCell[c=4]{l}\begin{enumerate}\item Professional CV, LinkedIn profile, cover letter, dan email profesional\item Technical writing dan dokumentasi proyek TI\item Presentasi teknis dan komunikasi tim multikultural\item Job interview: STAR method, behavioral dan technical interview\item Business English, IT portfolio, dan mock interview\end{enumerate} &  &  &  \\
\SetCell[r=2]{l} Pustaka & \SetCell[c=4]{l} \textbf{Utama:}\begin{enumerate}\item Mascull, B. 2017. \textit{Business Vocabulary in Use (Advanced)}. Cambridge.\item Guffey, M. 2019. \textit{Business Communication: Process and Product}. Cengage.\end{enumerate} &  &  &  \\
 & \SetCell[c=4]{l} \textbf{Pendukung:}\begin{enumerate}\item Swan, M. 2016. \textit{Practical English Usage}. Oxford.\end{enumerate} &  &  &  \\
Media Pembelajaran & \SetCell[c=2]{l} \textbf{Software:} Aplikasi presentasi, rekam audio/video, LMS, kamus daring. &  & \SetCell[c=2]{l} \textbf{Hardware:} Komputer/laptop, proyektor, mikrofon. &  \\
Nama Dosen Pengampu & \SetCell[c=4]{l}\begin{enumerate}\item Farid Angga Pribadi, S.Kom., M.Kom.\end{enumerate} &  &  &  \\
Matakuliah Syarat & \SetCell[c=4]{l} - &  &  &  \\
\end{longtblr}
\begin{longtblr}{width=\textwidth,colspec={|Q[c,m,0.06\textwidth]|X[1.35,l,m]|X[1.08,l,m]|X[1.16,l,m]|X[0.7,l,m]|X[1.24,l,m]|X[1.06,l,m]|X[0.98,l,m]|Q[c,m,0.05\textwidth]|},hlines,vlines,hspan=minimal,rowhead=2,row{1,2}={bg=headgreen,font=\bfseries},rowsep=1pt,colsep=1.5pt}
Mgg.\par Ke & Kemampuan Akhir Yang Direncanakan (Sub-CP-MK) & Bahan kajian (Materi Pembelajaran) & Bentuk dan Metode Pembelajaran & Estimasi Waktu & Pengalaman Belajar Mahasiswa & Kriteria \& Bentuk Penilaian & Indikator Penilaian & Bobot Penilaian (\%) \\
(1) & (2) & (3) & (4) & (5) & (6) & (7) & (8) & (9) \\
1 & SCPMK0303-05901, SCPMK0102-05902 & English for IT Professionals: overview dan kebutuhan pasar kerja global. & Modalitas: Blended Learning. Bentuk: Luring/Daring. Metode: Case Method, simulasi, role-play, dan peer feedback. & 1 x 2 x 50' tatap muka; 1 x 2 x 50' tugas/praktik mandiri. & Mahasiswa mengeksplorasi tren pasar kerja TI global, mengidentifikasi kompetensi Bahasa Inggris yang dibutuhkan, dan memetakan gap kompetensi diri. & Bentuk penilaian: Refleksi kebutuhan Bahasa Inggris dalam karir TI. Mengacu rubrik RTM. & Kelancaran komunikasi Bahasa Inggris; kualitas dokumen; profesionalisme. & 2 \\
2 & SCPMK0303-05901 & Professional CV dan LinkedIn profile dalam Bahasa Inggris. & Modalitas: Blended Learning. Bentuk: Luring/Daring. Metode: Case Method, simulasi, role-play, dan peer feedback. & 1 x 2 x 50' tatap muka; 1 x 2 x 50' tugas/praktik mandiri. & Mahasiswa menganalisis contoh CV dan LinkedIn profesional TI, merancang CV dan LinkedIn profile sendiri menggunakan konvensi profesional global. & Bentuk penilaian: Draft CV dan LinkedIn profile. Mengacu rubrik RTM. & Kelancaran komunikasi Bahasa Inggris; kualitas dokumen; profesionalisme. & 2 \\
3 & SCPMK0303-05901 & Cover letter dan application email: struktur dan konvensi. & Modalitas: Blended Learning. Bentuk: Luring/Daring. Metode: Case Method, simulasi, role-play, dan peer feedback. & 1 x 2 x 50' tatap muka; 1 x 2 x 50' tugas/praktik mandiri. & Mahasiswa mempelajari struktur cover letter dan email aplikasi kerja, kemudian menyusun cover letter dan email untuk posisi TI yang dipilih. & Bentuk penilaian: Draft cover letter dan application email. Mengacu rubrik RTM. & Kelancaran komunikasi Bahasa Inggris; kualitas dokumen; profesionalisme. & 2 \\
4 & SCPMK0303-05901 & Kuis 1: CV dan cover letter. & Modalitas: Blended Learning. Bentuk: Luring/Daring. Metode: Case Method, simulasi, role-play, dan peer feedback. & 1 x 2 x 50' tatap muka; 1 x 2 x 50' tugas/praktik mandiri. & Mahasiswa mengerjakan kuis tentang konvensi penulisan CV, cover letter, dan email profesional Bahasa Inggris dalam konteks TI. & Bentuk penilaian: Kuis 1 tertulis (CV dan cover letter). Mengacu rubrik RTM. & Kelancaran komunikasi Bahasa Inggris; kualitas dokumen; profesionalisme. & 5 \\
5 & SCPMK0303-05901 & Technical writing: dokumentasi dan laporan TI dalam Bahasa Inggris. & Modalitas: Blended Learning. Bentuk: Luring/Daring. Metode: Case Method, simulasi, role-play, dan peer feedback. & 1 x 2 x 50' tatap muka; 1 x 2 x 50' tugas/praktik mandiri. & Mahasiswa mempelajari prinsip technical writing dan menyusun dokumentasi teknis (README, API docs, atau laporan proyek) dalam Bahasa Inggris. & Bentuk penilaian: Sampel dokumentasi teknis TI. Mengacu rubrik RTM. & Kelancaran komunikasi Bahasa Inggris; kualitas dokumen; profesionalisme. & 2 \\
6 & SCPMK0102-05902 & Presentasi teknis: struktur, delivery, dan Q\&A dalam Bahasa Inggris. & Modalitas: Blended Learning. Bentuk: Luring/Daring. Metode: Case Method, simulasi, role-play, dan peer feedback. & 1 x 2 x 50' tatap muka; 1 x 2 x 50' tugas/praktik mandiri. & Mahasiswa berlatih menyusun dan menyampaikan presentasi teknis proyek TI dalam Bahasa Inggris, termasuk teknik menangani sesi tanya jawab. & Bentuk penilaian: Mini presentasi teknis (5 menit). Mengacu rubrik RTM. & Kelancaran komunikasi Bahasa Inggris; kualitas dokumen; profesionalisme. & 2 \\
7 & SCPMK0102-05902 & Komunikasi dalam tim multikultural: etika dan gaya komunikasi. & Modalitas: Blended Learning. Bentuk: Luring/Daring. Metode: Case Method, simulasi, role-play, dan peer feedback. & 1 x 2 x 50' tatap muka; 1 x 2 x 50' tugas/praktik mandiri. & Mahasiswa mempelajari perbedaan budaya dalam komunikasi profesional dan berlatih berkomunikasi dalam skenario tim TI multikultural. & Bentuk penilaian: Role-play komunikasi tim multikultural. Mengacu rubrik RTM. & Kelancaran komunikasi Bahasa Inggris; kualitas dokumen; profesionalisme. & 2 \\
8 & SCPMK0303-05901, SCPMK0102-05902 & UTS: portofolio dokumen profesional Bahasa Inggris. & Modalitas: Blended Learning. Bentuk: Luring/Daring. Metode: Case Method, simulasi, role-play, dan peer feedback. & 1 x 2 x 50' tatap muka; 1 x 2 x 50' tugas/praktik mandiri. & Mahasiswa mengumpulkan dan mempresentasikan portofolio dokumen profesional (CV, cover letter, email, technical writing) sebagai evaluasi tengah semester. & Bentuk penilaian: UTS portofolio (pengumpulan dan presentasi singkat). Mengacu rubrik RTM. & Kelancaran komunikasi Bahasa Inggris; kualitas dokumen; profesionalisme. & 15 \\
9 & SCPMK0102-05902 & Job interview preparation: common IT interview questions. & Modalitas: Blended Learning. Bentuk: Luring/Daring. Metode: Case Method, simulasi, role-play, dan peer feedback. & 1 x 2 x 50' tatap muka; 1 x 2 x 50' tugas/praktik mandiri. & Mahasiswa menganalisis pertanyaan interview umum di industri TI dan berlatih menjawab dengan bahasa Inggris yang profesional. & Bentuk penilaian: Latihan jawaban interview IT. Mengacu rubrik RTM. & Kelancaran komunikasi Bahasa Inggris; kualitas dokumen; profesionalisme. & 3 \\
10 & SCPMK0102-05902 & Behavioral interview: STAR method. & Modalitas: Blended Learning. Bentuk: Luring/Daring. Metode: Case Method, simulasi, role-play, dan peer feedback. & 1 x 2 x 50' tatap muka; 1 x 2 x 50' tugas/praktik mandiri. & Mahasiswa mempelajari dan berlatih STAR method (Situation, Task, Action, Result) untuk menjawab pertanyaan behavioral interview dalam Bahasa Inggris. & Bentuk penilaian: Latihan STAR method written dan lisan. Mengacu rubrik RTM. & Kelancaran komunikasi Bahasa Inggris; kualitas dokumen; profesionalisme. & 3 \\
11 & SCPMK0102-05902 & Kuis 2: interview skills dan workplace communication. & Modalitas: Blended Learning. Bentuk: Luring/Daring. Metode: Case Method, simulasi, role-play, dan peer feedback. & 1 x 2 x 50' tatap muka; 1 x 2 x 50' tugas/praktik mandiri. & Mahasiswa mengerjakan kuis tentang teknik interview kerja, STAR method, dan komunikasi profesional di tempat kerja. & Bentuk penilaian: Kuis 2 tertulis (interview skills). Mengacu rubrik RTM. & Kelancaran komunikasi Bahasa Inggris; kualitas dokumen; profesionalisme. & 5 \\
12 & SCPMK0102-05902 & Technical interview: coding discussion dan problem solving. & Modalitas: Blended Learning. Bentuk: Luring/Daring. Metode: Case Method, simulasi, role-play, dan peer feedback. & 1 x 2 x 50' tatap muka; 1 x 2 x 50' tugas/praktik mandiri. & Mahasiswa berlatih simulasi technical interview mencakup diskusi kode, problem solving, dan menjelaskan solusi teknis dalam Bahasa Inggris. & Bentuk penilaian: Simulasi technical interview singkat. Mengacu rubrik RTM. & Kelancaran komunikasi Bahasa Inggris; kualitas dokumen; profesionalisme. & 3 \\
13 & SCPMK0102-05902 & Business English: meeting, negotiation, dan email etiquette. & Modalitas: Blended Learning. Bentuk: Luring/Daring. Metode: Case Method, simulasi, role-play, dan peer feedback. & 1 x 2 x 50' tatap muka; 1 x 2 x 50' tugas/praktik mandiri. & Mahasiswa berlatih bahasa Inggris bisnis untuk konteks meeting profesional, negosiasi, dan penulisan email formal di lingkungan kerja TI. & Bentuk penilaian: Simulasi meeting dan email business. Mengacu rubrik RTM. & Kelancaran komunikasi Bahasa Inggris; kualitas dokumen; profesionalisme. & 3 \\
14 & SCPMK0303-05901 & IT portfolio: presentasi proyek dalam Bahasa Inggris. & Modalitas: Blended Learning. Bentuk: Luring/Daring. Metode: Case Method, simulasi, role-play, dan peer feedback. & 1 x 2 x 50' tatap muka; 1 x 2 x 50' tugas/praktik mandiri. & Mahasiswa menyusun dan berlatih mempresentasikan portofolio proyek TI dalam Bahasa Inggris sebagai persiapan UAS. & Bentuk penilaian: Draft presentasi portofolio proyek TI. Mengacu rubrik RTM. & Kelancaran komunikasi Bahasa Inggris; kualitas dokumen; profesionalisme. & 3 \\
15 & SCPMK0102-05902 & Mock interview: full simulation. & Modalitas: Blended Learning. Bentuk: Luring/Daring. Metode: Case Method, simulasi, role-play, dan peer feedback. & 1 x 2 x 50' tatap muka; 1 x 2 x 50' tugas/praktik mandiri. & Mahasiswa menjalani simulasi full mock interview (HR dan technical) dalam Bahasa Inggris dengan pasangan atau kelompok, mendapatkan feedback. & Bentuk penilaian: Rekaman simulasi mock interview. Mengacu rubrik RTM. & Kelancaran komunikasi Bahasa Inggris; kualitas dokumen; profesionalisme. & 3 \\
16 & SCPMK0102-05902, SCPMK0303-05901 & PBL: simulasi wawancara dan presentasi teknis final. & Modalitas: Blended Learning. Bentuk: Luring/Daring. Metode: Case Method, simulasi, role-play, dan peer feedback. & 1 x 2 x 50' tatap muka; 1 x 2 x 50' tugas/praktik mandiri. & Mahasiswa menjalani PBL berupa simulasi wawancara kerja komprehensif dan presentasi teknis portofolio proyek TI dalam Bahasa Inggris di depan instruktur. & Bentuk penilaian: PBL mock interview dan presentasi teknis (full assessment). Mengacu rubrik RTM. & Kelancaran komunikasi Bahasa Inggris; kualitas dokumen; profesionalisme. & 22 \\
17 & SCPMK0102-05902, SCPMK0303-05901 & UAS: final mock interview panel dan presentasi portofolio. & Modalitas: Blended Learning. Bentuk: Luring/Daring. Metode: Case Method, simulasi, role-play, dan peer feedback. & 1 x 2 x 50' tatap muka; 1 x 2 x 50' tugas/praktik mandiri. & Mahasiswa mengikuti UAS berupa mock interview panel (dihadapi tim instruktur) dan presentasi portofolio profesional final dalam Bahasa Inggris. & Bentuk penilaian: UAS final mock interview panel dan presentasi portofolio. Mengacu rubrik RTM. & Kelancaran komunikasi Bahasa Inggris; kualitas dokumen; profesionalisme. & 23 \\
\end{longtblr}
